% !TEX root = ../main.tex

\chapter{Ergebnisse}
\label{ch:ergebnisse}

\section{Deskriptive Auswertung}
\label{sec:deskriptiv}

\subsection{Zeitliche Entwicklung des Vertrauens in den Bundesrat}
\label{sec:zeitverlauf}

Die zentrale Fragestellung dieses Berichts betrifft die zeitliche Entwicklung des
Vertrauens in den Bundesrat, gemessen durch die Variable PP04 auf einer Skala von
0~(kein Vertrauen) bis 10~(volles Vertrauen). Tabelle~\ref{tab:trust-time} zeigt
die Mittelwerte und Mediane f\"ur alle 16~Erhebungsjahre.

\begin{table}[ht]
\centering
\caption{Mittleres und medianes Vertrauen in den Bundesrat nach Jahr}
\label{tab:trust-time}
\small
\begin{tabular}{@{}r r r @{\hspace{14pt}} r r r@{}}
\hline
\textbf{Jahr} & $\bar{x}$ & \textbf{Median} &
\textbf{Jahr} & $\bar{x}$ & \textbf{Median} \\
\hline
1999 & 5.74 & 6 & 2006 & 5.50 & 6 \\
2000 & 5.93 & 6 & 2007 & 5.51 & 6 \\
2001 & 5.92 & 6 & 2008 & 5.67 & 6 \\
2002 & 5.62 & 6 & 2009 & 5.39 & 5 \\
2003 & 5.44 & 6 & 2011 & 5.77 & 6 \\
2004 & 5.28 & 5 & 2014 & 5.98 & 6 \\
2005 & 5.39 & 5 & 2017 & 6.31 & 7 \\
     &      &   & 2020 & 6.76 & 7 \\
     &      &   & 2023 & 6.36 & 7 \\
\hline
\end{tabular}
\end{table}

\"Uber den gesamten Beobachtungszeitraum von 1999 bis 2023 zeigt sich ein insgesamt
steigender Trend. Zu Beginn der Erhebung im Jahr~1999 lag das mittlere Vertrauen bei
5.74 Skalenpunkten. In den folgenden Jahren sank es zun\"achst ab und erreichte im
Jahr~2004 mit 5.28 den tiefsten Wert des gesamten Beobachtungszeitraums. Dieser
R\"uckgang f\"allt zeitlich mit einer Phase zusammen, in der die Schweizer Politik durch
verschiedene kontroverse Themen gepr\"agt war, darunter Diskussionen \"uber die
bilateralen Vertr\"age mit der EU und innenpolitische Spannungen.

Ab 2005 setzt ein allm\"ahlicher Erholungstrend ein, der sich ab 2011 beschleunigt. Im
Jahr~2017 \"uberschreitet der Mittelwert erstmals die Marke von 6.0, und der Median
steigt von 6 auf 7. Den H\"ochstwert des gesamten Beobachtungszeitraums erreicht das
Vertrauen im Jahr~2020 mit einem Mittelwert von 6.76. Dieser markante Anstieg ist
zeitlich mit dem Ausbruch der COVID-19-Pandemie verbunden und l\"asst sich plausibel
als \emph{Rally-around-the-flag-Effekt} interpretieren: In Krisenzeiten tendieren
Bev\"olkerungen dazu, sich hinter der Regierung zu versammeln, was sich in einem
vor\"ubergehend erh\"ohten institutionellen Vertrauen \"aussert.

Im Jahr~2023 ist ein leichter R\"uckgang auf 6.36 zu beobachten, wobei der Wert
deutlich \"uber dem Niveau der fr\"uhen 2000er-Jahre liegt. Der Median bleibt bei~7,
was darauf hindeutet, dass die Mehrheit der Befragten weiterhin ein \"uberdurchschnittlich
hohes Vertrauen in den Bundesrat aufweist.

% TODO: Abbildung einbinden, sobald exportiert
%\begin{figure}[ht]
%    \centering
%    \includegraphics[width=0.85\textwidth]{Figures/trust_over_time.pdf}
%    \caption{Mittleres Vertrauen in den Bundesrat (PP04) \"uber die Zeit, 1999--2023}
%    \label{fig:trust-time}
%\end{figure}

\subsection{Vertrauen nach soziodemografischen Merkmalen}
\label{sec:untergruppen}

Um ein differenzierteres Bild der Determinanten des Vertrauens in den Bundesrat zu
erhalten, wurde die zeitliche Entwicklung von PP04 zus\"atzlich nach verschiedenen
soziodemografischen Untergruppen aufgeschl\"usselt. Die folgenden Abschnitte fassen
die wichtigsten Befunde zusammen.

\paragraph{Geschlecht.}
Die zeitlichen Verl\"aufe von M\"annern und Frauen verlaufen \"uber den gesamten
Beobachtungszeitraum nahezu parallel. Frauen weisen tendenziell ein leicht h\"oheres
mittleres Vertrauen auf als M\"anner, wobei die Differenz je nach Jahr zwischen 0.1 und
0.4~Skalenpunkten schwankt. Dieser Befund ist konsistent mit der Beobachtung, dass
Frauen in vielen europ\"aischen L\"andern ein leicht h\"oheres generalisiertes Vertrauen
in Institutionen berichten. Beide Geschlechter zeigen den gleichen Aufwerts- und
den gleichen COVID-bedingten Spitzenwert im Jahr~2020.

\paragraph{Geburtskohorte.}
Die Aufschl\"usselung nach Geburtskohorten zeigt ein deutliches Muster: \"Altere
Kohorten (insbesondere Jahrg\"ange vor 1950) weisen durchweg ein h\"oheres Vertrauen
auf als j\"ungere Kohorten. Die j\"ungste Kohorte (Jahrg\"ange 2000--2010), die erst ab
ca.\ 2018 in der Stichprobe vertreten ist, liegt bei ihrem Eintritt im unteren
Bereich. Trotz dieser Niveauunterschiede zeigen alle Kohorten den gleichen allgemeinen
Aufwertstrend. Die Kohortenunterschiede k\"onnten sowohl Alterseffekte (\"altere
Personen vertrauen generell mehr) als auch Kohorteneffekte (fr\"uher geborene
Generationen haben ein grundlegend anderes Verh\"altnis zu politischen Institutionen)
widerspiegeln -- eine Trennung dieser Effekte ist mit rein deskriptiven Mitteln
nicht m\"oglich.

\paragraph{Nationalit\"at.}
Personen mit ausschliesslich Schweizer Nationalit\"at (CH) und Personen mit doppelter
Staatsb\"urgerschaft (CH+Andere) unterscheiden sich im mittleren Vertrauen nur
geringf\"ugig. Personen mit ausschliesslich ausl\"andischer Nationalit\"at zeigen ein
leicht h\"oheres Vertrauen, wobei diese Gruppe deutlich kleiner ist und die Sch\"atzungen
entsprechend unsicherer sind. Eine m\"ogliche Erkl\"arung ist, dass ausl\"andische
Staatsangeh\"orige die Schweizer Institutionen im Vergleich zu ihrem Herkunftsland
positiver bewerten.

\paragraph{Eigene Kinder.}
Der Vergleich zwischen Personen mit und ohne eigene Kinder zeigt nur geringe
Unterschiede. Personen mit Kindern weisen ein tendenziell leicht h\"oheres Vertrauen
auf, was m\"oglicherweise mit einer st\"arkeren Einbindung in gesellschaftliche
Strukturen (Schule, Vereine, Nachbarschaft) zusammenh\"angt.

\paragraph{Gesundheitliche Einschr\"ankung.}
Eines der deutlichsten Muster zeigt sich bei der Aufschl\"usselung nach der Variable
PC19 (langfristige gesundheitliche Einschr\"ankung). Personen mit einer Einschr\"ankung
berichten ein deutlich niedrigeres Vertrauen in den Bundesrat als Personen ohne
Einschr\"ankung. Die Differenz betr\"agt in einigen Jahren \"uber einen Skalenpunkt und
ist \"uber den gesamten Beobachtungszeitraum stabil. Dieser Befund k\"onnte darauf
hindeuten, dass Personen mit gesundheitlichen Problemen h\"aufiger negative Erfahrungen
mit staatlichen Leistungen und Versicherungssystemen machen, was sich in einem
geringeren Institutionenvertrauen niederschl\"agt.

\paragraph{Zivilstand.}
Die Vertrauenswerte unterscheiden sich nach Zivilstand, wobei die Muster im
folgenden Abschnitt ausf\"uhrlich diskutiert werden. Vorab sei erw\"ahnt, dass die
h\"ochsten Werte bei Personen in eingetragener Partnerschaft und die niedrigsten bei
geschiedenen Personen zu beobachten sind. Im zeitlichen Verlauf zeigen alle
Zivilstandsgruppen einen \"ahnlichen Aufwertstrend, wobei die Abst\"ande zwischen den
Gruppen weitgehend stabil bleiben.

\section{Auswertung der Hypothese}
\label{sec:hypothese-auswertung}

Die in Abschnitt~\ref{sec:hypothese} formulierte Hypothese besagt, dass verheiratete
Personen ein h\"oheres Vertrauen in den Bundesrat aufweisen als ledige, getrennte oder
geschiedene Personen. Tabelle~\ref{tab:civsta} zeigt die deskriptiven Kennwerte
nach Zivilstand, sortiert nach absteigendem Mittelwert.

\begin{table}[ht]
\centering
\caption{Vertrauen in den Bundesrat nach Zivilstand}
\label{tab:civsta}
\begin{tabular}{@{}l r r r@{}}
\hline
\textbf{Zivilstand} & $\bar{x}$ \textbf{PP04} & \textbf{Median} & \textbf{n} \\
\hline
Eingetragene Partnerschaft & 6.51 & 7 & 218 \\
Aufgel\"oste Partnerschaft   & 6.36 & 6 & 11 \\
Verheiratet                & 5.96 & 6 & 95\,195 \\
Verwitwet                  & 5.94 & 6 & 7\,384 \\
Ledig                      & 5.86 & 6 & 45\,062 \\
Getrennt                   & 5.61 & 6 & 2\,145 \\
Geschieden                 & 5.49 & 6 & 12\,488 \\
\hline
\end{tabular}
\end{table}

Die Ergebnisse lassen sich wie folgt zusammenfassen. Verheiratete Personen weisen mit
einem Mittelwert von 5.96 tats\"achlich ein h\"oheres mittleres Vertrauen auf als ledige
(5.86), getrennte (5.61) und geschiedene Personen (5.49). Die Richtung des Zusammenhangs
entspricht somit der formulierten Hypothese.

Allerdings sind die Unterschiede unterschiedlich stark ausgepr\"agt. Die Differenz
zwischen Verheirateten und Ledigen betr\"agt lediglich 0.10~Skalenpunkte und ist damit
gering. Deutlicher ist die Differenz zu geschiedenen Personen mit 0.47~Skalenpunkten
und zu getrennten Personen mit 0.35~Skalenpunkten. Dies deutet darauf hin, dass nicht
die Ehe an sich, sondern vielmehr die Erfahrung einer Trennung oder Scheidung mit
einem niedrigeren Vertrauen assoziiert ist.

Bemerkenswert ist der hohe Mittelwert bei Personen in eingetragener Partnerschaft
(6.51), der alle anderen Gruppen \"ubertrifft. Allerdings ist die Fallzahl dieser
Gruppe mit $n = 218$ vergleichsweise gering, sodass dieser Wert mit Vorsicht zu
interpretieren ist. Auch die Kategorie ``aufgel\"oste Partnerschaft'' ($n = 11$) ist
aufgrund der minimalen Fallzahl statistisch nicht belastbar.

Im zeitlichen Verlauf zeigen alle Zivilstandsgruppen den bereits beschriebenen
allgemeinen Aufwertstrend. Die Rangfolge der Gruppen bleibt dabei \"uber die Jahre
hinweg weitgehend stabil: Geschiedene und getrennte Personen liegen durchweg unter
dem Gesamtdurchschnitt, w\"ahrend Verheiratete leicht dar\"uber liegen.

Zusammenfassend l\"asst sich die Hypothese \textbf{teilweise best\"atigen}. Der
Zivilstand ist mit dem Vertrauen in den Bundesrat assoziiert, wobei die Richtung des
Zusammenhangs der Erwartung entspricht. Verheiratete Personen zeigen ein h\"oheres
Vertrauen als geschiedene und getrennte Personen. Die Differenz zu ledigen Personen
ist jedoch gering und deutet darauf hin, dass weniger der Familienstand an sich als
vielmehr die Erfahrung eines Beziehungsabbruchs mit einem niedrigeren Vertrauen
einhergeht. Eine kausale Interpretation ist auf Grundlage der rein deskriptiven
Analyse nicht m\"oglich, da Drittvariablen wie Alter, Einkommen und Bildung nicht
kontrolliert wurden.
